\documentclass[12pt]{article}
\usepackage{anysize}
\usepackage{xltxtra}
\usepackage{xunicode}
\usepackage{graphicx}
\usepackage{color}
\usepackage{tikz}
\usepackage{minted}
\usepackage{listings}
\usepackage{enumitem}
\usepackage{framed}
\usepackage{relsize}
\usepackage{float}
\usepackage{algpseudocode}
\usepackage[hyphens]{url}

\usepackage{color,hyperref}

\usepackage{fancyhdr}
\usepackage{lastpage}
\usepackage{fontspec}

\linespread{1.5}

\definecolor{darkblue}{rgb}{0.0,0.0,0.3}
\hypersetup{colorlinks,breaklinks,
        linkcolor=darkblue,urlcolor=darkblue,
            anchorcolor=black,citecolor=black}

\def\name{Grigorios Lyras}
\def\IF{\textbf{if}}
\def\THEN{\textbf{then}}
\def\SAYS{\textbf{says}}
\def\SPEAKSFOR{$\Rightarrow$}
\def\FOR{\textbf{for}}
\def\AND{\textbf{and}}

\setlength{\headheight}{15.2pt}

\pagestyle{fancy}{ %
    \fancyhf{} % remove everything
    \renewcommand{\headrulewidth}{0pt} % remove lines as well
    \renewcommand{\footrulewidth}{0.5pt}
    \lfoot{\name}
    \rfoot{Page \thepage\ of \pageref{LastPage}}

\renewcommand{\labelenumi}{(\alph{enumi})}
\defaultfontfeatures{Ligatures=TeX} % To support LaTeX quoting style
\setmainfont[Mapping=TeX-text]{CMU Serif}


\begin{document}

\begin{titlepage}

\begin{tikzpicture}[remember picture, overlay]
      \node [anchor=north east, inner sep=0pt]  at (current page.north east)
           {\includegraphics[height=3cm]{title/ucl_logo}};
       \end{tikzpicture}
\begin{center}
\begin{LARGE}\textbf{
University College London\\}\end{LARGE}
\vspace{1cm}
\begin{Large}Simulating Path Planning for Data\end{Large}\\
\begin{Large}Muling in Wireless Sensor Networks\end{Large}\\
\begin{Large}using Unmanned Aerial Vehicles\end{Large}\\
\vspace{2cm}
\begin{LARGE}
MSc Networked Computer Systems\\
NCS Group Project Individual Report\\
\end{LARGE}
\large{Academic Year 2013-2014}\\
\vfill
\end{center}
This report is submitted as part requirement for the MSc in Networked Computer
Systems at University College London. It is substantially the result of our own
work except where explicitly indicated in the text. The report may be freely
copied and distributed provided the source is explicitly acknowledged.\\
\begin{center}
\begin{tabular}{l r}
\Large{Grigorios Lyras}&
\large{SID: 13047325}\\
\end{tabular}\\
\vspace{1cm}
\large\today\\
\end{center}
\end{titlepage}



\tableofcontents{}

\pagebreak

\section{Introduction}
The aim of this report is the discussion of the Group Project that I took part
in as part of the NCS team, comprised of myself and Seng Jea Lee, Johann Mifsud,
Minmin Zhang. In this report I will present my own role in the progress of
the project. The analysis will be broken into three main sections. In objectives
and achievements I will present the objectives we set through the duration of
the project as well as the achievements and my contribution to them. Moving to
the results of the project we will discuss the output of the project as well as
the quality of the outcome. In the third section, I will present my view on the
project management techniques and tools that we used along with my experience on
them. The interaction with the rest of the team members and the task
distribution will also be reported. In the final section I will present my own
contribution to the project.

\section{Project Overview}
Wireless Sensor Networks are systems of sensors distributed over an area in
order to collect useful information. These are low cost devices with limitations
in their capabilities. One of these limitations is the communication range. As a
result it is understood that these networks are not always connected. A concept
to overcome such difficulties is the use of data mules that visit these devices
and collect the data from them. These data mules are networked devices that are
able to move from location to location, pick up the data, store them
temporarily and then deliver them back to the base. An interesting area is the
exploration of data muling with the use of Unmanned Aerial Vehicles (UAVs).
UAVs have unique mobility abilities which make them a very suitable solution
for data muling.

Since the suitability of UAVs for the task is established we move to the
implementation. In order to ensure that UAVs behave in a safe manner thorough
testing is required. However this is not always possible due to weather
conditions. Another factor that increases the risk significantly in this
research area is the cost of the UAVs the nearby environment. The high cost
makes field testing unsuitable for this research area. As a consequence the
research has focused on simulating UAVs in order to reduce the risk mentioned
above.

The initial idea for this project was the exploration of path planning
algorithms for UAVs when they are used to collect data from distributed wireless
sensor networks. Due to the limitations and risk involved in such a project if it
is taken on the field the exploration was limited to a simulated environment.
After researching the area we focused on a simulator that enables the physical
simulation of UAVs. CRATES was the simulator we used for this project which
lacked support for modeling the wireless channel.

We identified the areas where the project would contribute to. This project set
to build a radio propagation model for CRATES and utilise this in order to
explore higher level Path Planning Algorithms. These were to be evaluated on
their performance based on distance and data collection statistics.


\section{Objectives and Achievements}
One of the most challenging aspects of this project has been it's wide scope. It
is in effect a multi disciplinary project. The main objectives were the
extension of CRATES in order to support wireless communication as well as
the implementation and testing of path planning algorithms. However, these two
main objectives were divided into multiple milestones along the progress of the
project.

\subsection{Wireless Channel Modeling}
Taking into account the project outline one of the first milestones along the
timeline was the Modeling of the wireless channel in CRATES. This task was split
into two main subgoals. Firstly we had to comprehend how CRATES operated and how
it would be best to implement the desired functionality. In parallel we were to
research on how to model the wireless channel.

CRATES supports wireless sensors however, it is limited because it offers no
communication model. In order to add support for the communication model we
extended CRATES and enabled the models inside the simulated world to pass
messages from one to another. These communication primitives were to be enriched
with a realistic radio model that would control the radio propagation method.
Since we required a detailed wireless communication simulation, it was
essential to simulate not only the raw wireless channel, but also the network
stack layers up to the internet protocol.

In order to simulate the network stack layers we reused existing code. In
particular we used ns2, a discrete event simulator, that executes TCL scripts at
runtime. ns2 is a widely used simulator in research areas and supports multiple
MAC layers. Alternatives to ns2 were considered, however, due to our previous
familiarity with ns2 was what led us to use ns2 for our project. In our
implementation we run two simulators in parallel. We treat the ns2 simulation
as a function call, pushing the input, waiting for a fixed amount of time and
then parsing the result.

In the area of wireless modeling we had to gather data for the network modules
that were available on the real life platforms. To do that we created a
measurement platform using OrisenStamps and NVO8C GPS units that allowed us to
measure wireless data as well as GPS coordinates. This way we could calculate
the distances between the different platforms. These platforms were flown in the
air by two UAVs in order to gather the required data. The OrisenStamps are
equipped with ZigBee 802.15.4 transceivers which we used to gather the ZigBee
experiment data set. We also gathered data for the WiFi channel. For the WiFi
measurements we used the onboard 802.11g WiFi modules that are onboard the
Fit-PC2 platforms which were also attached to the UAVs.

\subsection{Path Planning}
The second main goal for our project has been the exploration of path planning
algorithms. We focused on two main algorithms; one based on Sugihara's path
planning algorithm and a modified genetic algorithm. These two algorithms are
tested and their performance is compared on the basis of total distance
travelled by the UAVs and the amount of data collected.

\section{Project Results}
In this section of the report we will discuss the output of the project on the
goals set in radio modeling, the simulator plugin and the path planning
algorithms.

\subsection{Radio Modeling}
In the area of radio modeling we set a goal for accurately simulating the
wireless channel. In order to achieve this goal, initially we gathered wireless
communication data for two different MAC protocols, ZigBee 802.15.4 and WiFi
802.11g. For the ZigBee measurements we used a pair of OrisenStamps which exchanged
packets. From these packets we extracted the receiver signal strength
indication. In a similar manner we collected data for the WiFi network card that
was onboard to the Fit-PC2 that was used in our experiments. We observed a
fairly expected distribution of RSSI measurements over distance for the ZigBee
modules. The measurements taken showed an unrealistic path loss exponent of 1.07
whereas for the WiFi the exponent was about 1.41. These values are most likely
attributed to the unknown methods of the calculation of the RSSI values both by
the ZigBee and the WiFi modules. Moreover, for the WiFi measurements, another
possible explanation would be the limitations that were imposed on the allowed
flight distance and the coding scheme changes of the 802.11g MAC layer.
\subsection{Simulator Plugin}
The main software component of this project was the simulator plugin. This
enables the simulation of the wireless channel in CRATES. It provides the
required functionality and when compared with the real life measurements we
gathered it showed minimal discrepancy.
\subsection{Path Planning}
The final milestone of the project was the exploration of the two path planning
algorithms. Both Sugihara's algorithm and the optimized genetic algorithm were
implemented and compared on the basis of distance travelled by the UAVs. It can
be seen as in the group report that the optimized algorithm saves more distance
as the number of nodes increases.

\section{Project Management}
An important part of the project has been the project management aspect. Since
this was an exploration project with multiple milestones we required flexibility
in the goals that we set for every week. We kept the meetings informal where we
exchanged information so that everyone is always up to date. Moreover we aimed
to work often in the same space to keep everyone motivated and our productivity
high. When that was not possible we tried to have a constant online presence
that enabled our group to exchange ideas and troubleshoot in a very agile
manner.

Throughout the project we had a very versatile and flexible organisation that
enabled members to pick up tasks that were more fitting to their interest and
skillset. In order to keep track we used a number of tools. One of the main
techniques we applied was the use of the Human Changelog where all tasks and
their progress were recorded for everyone to see and update. In addition we used
Trello to manage our tasks, which seemed promising in the beginning however it
proved a burden instead because we were already using the Human Changelog.

An important part of this project has been software development. To be exact it
has been a choir of multiple smaller software projects. In order to cope with
the management of the codebase we used a distributed version control system that
allowed everyone to contribute to the task at hand. This was very helpful,
especially when we were debugging as we were able to push and pull changes
between our members with ease.


\subsection{Team Management}
This project took place in a highly dynamic team. We established our interests
since the beginning of the project and these were the main compass on how the
tasks were distributed to the members. The role of the Project Manager was
assigned to Seng Jea who was responsible for the team schedule as well as the
meetings with Prof. Hailes.

The team dynamics enabled us to pursue such a wide scoped project. That being
said specific limitations of our team management strategy crept in during our
project. One of these limitations we faced, was what resulted in this project
becoming so widely scoped. For example, the fact that we did not keep a stricter
Changelog turned out to hinder the evaluation of our progress during this
project. Moreover, for the same reason, we did not have deliverables to present
on a weekly basis as the task fragmentation was rather coarse grained. Overall
the team worked well as a whole, each member using their strengths and
complementing other members skills.

\subsubsection{Seng Jea Lee}
Seng Jea was assigned the role of the Project Manager. He has showed deep
analytical skills and profound understanding of the field of wireless modeling.
He shows extreme attention to detail and professionalism. In his role as a
Project Manager he kept everyone motivated and in line with the requirements of
the task at hand.

\subsubsection{Johann Mifsud}
Johann is an experienced Software Engineer with a strong background. He
has produced high quality deliverables adhering to deadlines set by the team.
He is hard working and dependable.

\subsubsection{Minmin Zhang}
Minmin is a very knowledgeable and very hard working individual. Her work on the
research area of the project was of remarkable quality and depth. Moreover she
did not set aside the task of implementing the path planning algorithms despite
their complexity. That being said it would seem that Minmin was often shy to
ask for help during the project.

In the long run, in my opinion the tasks were fairly evenly distributed among
the members of the team.

\subsubsection{Personal Workplan}
Throughout this project I engaged with a number of diverse tasks. In the
beginning of the project I enthusiastically began the setup of the simulation
framework that we were assigned to develop upon. Following that, I engaged with
setting up a constant integration server for our project, at which I was
unsuccessful due to specific requirements of the framework not being available
within online services. Subsequently I focused on understanding the platforms
that we were planning to use for our experiments. I was assigned with the task
of understanding what information were we able to extract from them and identify
the optimal method to gather these data. This task proved to be hindered by the
fact that the manual for the Fit-PC2 platform was unavailable to us. Finally, I
was assigned with the task of gathering the data, a process that involved fully
understanding the platforms, resolving issues such as failed chips and multiple
trips to RVC. During the final phase of the project I was assigned with my share
of the Group Report. Due to some unexpected bugs to our CRATES plugin, I also
undertook the task of debugging the plugin.

Overall, I consider my contribution was commensurate to my abilities and has
positively complemented the efforts put by the other members in this project. I
was assigned many diverse tasks that provided me a complete overview of the
project and stimulated by attention and enthusiasm.
\end{document}
